\PageTitle{05 / Run variation}{Batch-level latency distribution}
The raw evidence retains batch membership. Figure~\ref{fig:batch} shows each batch's nearest-rank median and minimum--maximum range. It deliberately uses a batch index, not a synthetic wall-clock timeline.

\begin{figure}[h]\centering\begin{tikzpicture}
\begin{axis}[xmin=.5,xmax=10.5,ymin=0,ymax=750,xtick={1,...,10},xlabel={Recorded batch index (not elapsed time)},ylabel={Client-observed latency (ms)},ymajorgrids=true,legend pos=north east,height=65mm]
\addplot[MutedInk,line width=1.5pt,forget plot] coordinates {(1,628) (1,695)};
\addplot[MutedInk,line width=1.5pt,forget plot] coordinates {(2,205) (2,278)};
\addplot[MutedInk,line width=1.5pt,forget plot] coordinates {(3,214) (3,278)};
\addplot[MutedInk,line width=1.5pt,forget plot] coordinates {(4,465) (4,619)};
\addplot[MutedInk,line width=1.5pt,forget plot] coordinates {(5,316) (5,450)};
\addplot[MutedInk,line width=1.5pt,forget plot] coordinates {(6,351) (6,416)};
\addplot[MutedInk,line width=1.5pt,forget plot] coordinates {(7,222) (7,292)};
\addplot[MutedInk,line width=1.5pt,forget plot] coordinates {(8,342) (8,501)};
\addplot[MutedInk,line width=1.5pt,forget plot] coordinates {(9,354) (9,411)};
\addplot[MutedInk,line width=1.5pt,forget plot] coordinates {(10,302) (10,372)};
\addplot[ConversationBlue,only marks,mark=*,mark size=2.2pt] coordinates {(1,662) (2,246) (3,245) (4,500) (5,345) (6,384) (7,256) (8,377) (9,385) (10,340)};\addlegendentry{Median; vertical line = min--max}
\end{axis}
\end{tikzpicture}
\caption{Recorded latency range and median for each ten-call batch. The first batch includes connection warm-up according to the source report.}\label{fig:batch}\end{figure}
\Source{S1 and S2; batch summaries recomputed from raw samples}

\subsection*{Observed variation}
The first batch spans 628--695 ms, with a median of 662 ms. Batches 2 and 3 have medians of 246 and 245 ms. Batch 4 rises to 500 ms. These differences show variation within the short run; the available fields do not establish its causes.

The source report states that the first batch includes connection warm-up. All 100 samples remain in the overall analysis. Removing warm-up observations after viewing the results would define a different analysis and must be labelled separately.

\begin{table}[h]\caption{Interpretation of workload and rate terms}\EvidenceTableStart
\begin{tabularx}{\linewidth}{@{}P{38mm}X@{}}\toprule
Term & Meaning in this report\\\midrule
Burst batch size & Ten calls launched concurrently, as reported by the source\\
Batch interval & Reported start pacing of at least one second\\
Completion rate & 100 completions divided by 11.946 seconds of total wall time\\
Maximum sustainable rate & Not measured; no saturation search or steady-state campaign\\
Resource efficiency & Not measured; CPU, memory, and database metrics absent\\\bottomrule
\end{tabularx}\end{table}

\Limit{Do not present this as an evenly spaced 10-RPS stream or a sustained 10-RPS capacity test. The JSON cannot verify overlap, queue depth, server utilization, or latency experienced by other users.}
